
\documentclass{article}
\usepackage{colortbl}
\usepackage{makecell}
\usepackage{multirow}
\usepackage{supertabular}

\begin{document}

\newcounter{utterance}

\centering \large Interaction Transcript for game `bbh', experiment `formal\_fallacies', episode 160 with selfrefine{-}Qwen3.5{-}27B.
\vspace{24pt}

{ \footnotesize  \setcounter{utterance}{1}
\setlength{\tabcolsep}{0pt}
\begin{supertabular}{c@{$\;$}|p{.15\linewidth}@{}p{.15\linewidth}p{.15\linewidth}p{.15\linewidth}p{.15\linewidth}p{.15\linewidth}}
    \# & \multicolumn{2}{c}{Player} && \multicolumn{2}{c}{Game Master} \\
    \hline

    \theutterance \stepcounter{utterance}  
    & & & \multicolumn{4}{p{0.6\linewidth}}{
        \cellcolor[rgb]{0.9,0.9,0.9}{
            \makecell[{{p{\linewidth}}}]{
                \texttt{\tiny{[P1$\langle$GM]}}
                \texttt{Distinguish deductively valid arguments from formal fallacies.} \\
\\ 
\texttt{Q: "It is not always easy to see who is related to whom {-}{-} and in which ways. The following argument pertains to this question: To begin with, Lesley is a close friend of Fernando. Moreover, being a close friend of Fernando or a schoolmate of Lowell is sufficient for being a great{-}grandfather of Leroy. It follows that Lesley is a great{-}grandfather of Leroy."} \\
\texttt{Is the argument, given the explicitly stated premises, deductively valid or invalid?} \\
\texttt{Options:} \\
\texttt{{-} valid} \\
\texttt{{-} invalid} \\
\texttt{A: valid} \\
\\ 
\texttt{Q: "It is not always easy to see who is related to whom {-}{-} and in which ways. The following argument pertains to this question: Whoever is not a great{-}grandfather of Clyde is a stepbrother of Brian. Being an ancestor of Dana is sufficient for not being a great{-}grandfather of Clyde. We may conclude: Everyone who is an ancestor of Dana is a stepbrother of Brian, too."} \\
\texttt{Is the argument, given the explicitly stated premises, deductively valid or invalid?} \\
\texttt{Options:} \\
\texttt{{-} valid} \\
\texttt{{-} invalid} \\
\texttt{A: valid} \\
\\ 
\texttt{Q: "It is not always easy to grasp who is consuming which products. The following argument pertains to this question: Every infrequent user of Paul Mitchell shampoo is either a rare consumer of Nioxin shampoo or a loyal buyer of Caress soap, or both. No regular consumer of Lush soap is a rare consumer of Nioxin shampoo and, in the same time, a loyal buyer of Caress soap. It follows that whoever is an infrequent user of Paul Mitchell shampoo is not a regular consumer of Lush soap."} \\
\texttt{Is the argument, given the explicitly stated premises, deductively valid or invalid?} \\
\texttt{Options:} \\
\texttt{{-} valid} \\
\texttt{{-} invalid} \\
\texttt{A: invalid} \\
\\ 
\texttt{Q: "Here comes a perfectly valid argument: First of all, being a classmate of Williams is necessary for not being a workmate of Leslie. Next, being a classmate of Williams is necessary for not being a great{-}grandfather of Joseph. Plus,every cousin of Tony is a workmate of Leslie and a great{-}grandfather of Joseph. Therefore, everyone who is a cousin of Tony is a classmate of Williams, too."} \\
\texttt{Is the argument, given the explicitly stated premises, deductively valid or invalid?} \\
\texttt{Options:} \\
\texttt{{-} valid } \\
\texttt{{-} invalid} \\
\texttt{A:} \\
            }
        }
    }
     \\ \\

    \theutterance \stepcounter{utterance}  
    & & \multicolumn{4}{p{0.6\linewidth}}{
        \cellcolor[rgb]{0.95,0.95,0.95}{
            \makecell[{{p{\linewidth}}}]{
                \texttt{\tiny{[P1$\leftrightarrow$P]}}
                \texttt{{[}DRAFT{]}} \\
\texttt{To determine the validity of the argument, let's break it down into logical propositions and analyze the structure step by step.} \\
\\ 
\texttt{**1. Define the Propositions:**} \\
\texttt{Let \$C\$ = is a classmate of Williams.} \\
\texttt{Let \$W\$ = is a workmate of Leslie.} \\
\texttt{Let \$G\$ = is a great{-}grandfather of Joseph.} \\
\texttt{Let \$T\$ = is a cousin of Tony.} \\
\\ 
\texttt{**2. Translate the Premises:**} \\
\\ 
\texttt{*   **Premise 1:** "Being a classmate of Williams is necessary for not being a workmate of Leslie."} \\
\texttt{    *   Logical form: If not \$W\$, then \$C\$. (\$\textbackslash{}neg W \textbackslash{}rightarrow C\$)} \\
\texttt{    *   Contrapositive: If not \$C\$, then \$W\$. (\$\textbackslash{}neg C \textbackslash{}rightarrow W\$)} \\
\\ 
\texttt{*   **Premise 2:** "Being a classmate of Williams is necessary for not being a great{-}grandfather of Joseph."} \\
\texttt{    *   Logical form: If not \$G\$, then \$C\$. (\$\textbackslash{}neg G \textbackslash{}rightarrow C\$)} \\
\texttt{    *   Contrapositive: If not \$C\$, then \$G\$. (\$\textbackslash{}neg C \textbackslash{}rightarrow G\$)} \\
\\ 
\texttt{*   **Premise 3:** "Every cousin of Tony is a workmate of Leslie and a great{-}grandfather of Joseph."} \\
\texttt{    *   Logical form: If \$T\$, then (\$W\$ and \$G\$). (\$T \textbackslash{}rightarrow (W \textbackslash{}land G)\$)} \\
\\ 
\texttt{**3. Analyze the Conclusion:**} \\
\texttt{*   **Conclusion:** "Everyone who is a cousin of Tony is a classmate of Williams, too."} \\
\texttt{    *   Logical form: If \$T\$, then \$C\$. (\$T \textbackslash{}rightarrow C\$)} \\
\\ 
\texttt{**4. Evaluate the Deduction:**} \\
\texttt{We need to see if the conclusion (\$T \textbackslash{}rightarrow C\$) necessarily follows from the premises.} \\
\\ 
\texttt{*   From Premise 3, we know that if someone is a cousin of Tony (\$T\$), they are both a workmate of Leslie (\$W\$) and a great{-}grandfather of Joseph (\$G\$).} \\
\texttt{*   So, for any person \$x\$ who is a cousin of Tony, we have \$W(x)\$ is true and \$G(x)\$ is true.} \\
\texttt{*   Now look at Premise 1 and Premise 2. They tell us what happens if someone is **not** a workmate (\$\textbackslash{}neg W\$) or **not** a great{-}grandfather (\$\textbackslash{}neg G\$). Specifically, they state that \$\textbackslash{}neg W \textbackslash{}rightarrow C\$ and \$\textbackslash{}neg G \textbackslash{}rightarrow C\$.} \\
\texttt{*   However, the premises do **not** state what happens if someone *is* a workmate (\$W\$) or *is* a great{-}grandfather (\$G\$).} \\
\texttt{    *   Knowing \$W\$ is true does not allow us to infer \$C} \\
            }
        }
    }
    & \\ \\

    \theutterance \stepcounter{utterance}  
    & & \multicolumn{4}{p{0.6\linewidth}}{
        \cellcolor[rgb]{0.95,0.95,0.95}{
            \makecell[{{p{\linewidth}}}]{
                \texttt{\tiny{[P1$\leftrightarrow$P]}}
                \texttt{{[}SELF{-}FEEDBACK{]}} \\
\texttt{FORMAT: problem} \\
\texttt{RULES: problem} \\
\texttt{STRATEGY: ok} \\
\texttt{REASON: The answer draft fails to provide the required final selection ("valid" or "invalid") and is cut off mid{-}sentence, violating the game's requirement for a complete answer.} \\
            }
        }
    }
    & \\ \\

    \theutterance \stepcounter{utterance}  
    & \multicolumn{4}{p{0.6\linewidth}}{
        \cellcolor[rgb]{0.561,0.737,0.561}{
            \makecell[{{p{\linewidth}}}]{
                \texttt{\tiny{[P1$\rangle$GM]}}
                \texttt{To determine the validity of the argument, let's break it down into logical propositions and analyze the structure step by step.} \\
\\ 
\texttt{**1. Define the Propositions:**} \\
\texttt{Let \$C\$ = is a classmate of Williams.} \\
\texttt{Let \$W\$ = is a workmate of Leslie.} \\
\texttt{Let \$G\$ = is a great{-}grandfather of Joseph.} \\
\texttt{Let \$T\$ = is a cousin of Tony.} \\
\\ 
\texttt{**2. Translate the Premises:**} \\
\\ 
\texttt{*   **Premise 1:** "Being a classmate of Williams is necessary for not being a workmate of Leslie."} \\
\texttt{    *   Logical form: If not \$W\$, then \$C\$. (\$\textbackslash{}neg W \textbackslash{}rightarrow C\$)} \\
\texttt{    *   Contrapositive: If not \$C\$, then \$W\$. (\$\textbackslash{}neg C \textbackslash{}rightarrow W\$)} \\
\\ 
\texttt{*   **Premise 2:** "Being a classmate of Williams is necessary for not being a great{-}grandfather of Joseph."} \\
\texttt{    *   Logical form: If not \$G\$, then \$C\$. (\$\textbackslash{}neg G \textbackslash{}rightarrow C\$)} \\
\texttt{    *   Contrapositive: If not \$C\$, then \$G\$. (\$\textbackslash{}neg C \textbackslash{}rightarrow G\$)} \\
\\ 
\texttt{*   **Premise 3:** "Every cousin of Tony is a workmate of Leslie and a great{-}grandfather of Joseph."} \\
\texttt{    *   Logical form: If \$T\$, then (\$W\$ and \$G\$). (\$T \textbackslash{}rightarrow (W \textbackslash{}land G)\$)} \\
\\ 
\texttt{**3. Analyze the Conclusion:**} \\
\texttt{*   **Conclusion:** "Everyone who is a cousin of Tony is a classmate of Williams, too."} \\
\texttt{    *   Logical form: If \$T\$, then \$C\$. (\$T \textbackslash{}rightarrow C\$)} \\
\\ 
\texttt{**4. Evaluate the Deduction:**} \\
\texttt{We need to see if the conclusion (\$T \textbackslash{}rightarrow C\$) necessarily follows from the premises.} \\
\\ 
\texttt{*   From Premise 3, we know that if someone is a cousin of Tony (\$T\$), they are both a workmate of Leslie (\$W\$) and a great{-}grandfather of Joseph (\$G\$).} \\
\texttt{*   So, for any person \$x\$ who is a cousin of Tony, we have \$W(x)\$ is true and \$G(x)\$ is true.} \\
\texttt{*   Now look at Premise 1 and Premise 2. They tell us what happens if someone is **not** a workmate (\$\textbackslash{}neg W\$) or **not** a great{-}grandfather (\$\textbackslash{}neg G\$). Specifically, they state that \$\textbackslash{}neg W \textbackslash{}rightarrow C\$ and \$\textbackslash{}neg G \textbackslash{}rightarrow C\$.} \\
\texttt{*   However, the premises do **not** state what happens if someone *is* a workmate (\$W\$) or *is* a great{-}grandfather (\$G\$).} \\
\texttt{    *   Knowing \$W\$ is true does not allow us to infer \$C} \\
            }
        }
    }
    & & \\ \\

    \theutterance \stepcounter{utterance}  
    & & & \multicolumn{2}{p{0.3\linewidth}}{
        \cellcolor[rgb]{0.9,0.9,0.9}{
            \makecell[{{p{\linewidth}}}]{
                \texttt{\tiny{[GM$|$GM]}}
                \texttt{Target: invalid} \\
            }
        }
    }
    & & \\ \\

    \theutterance \stepcounter{utterance}  
    & & & \multicolumn{2}{p{0.3\linewidth}}{
        \cellcolor[rgb]{0.9,0.9,0.9}{
            \makecell[{{p{\linewidth}}}]{
                \texttt{\tiny{[GM$|$GM]}}
                \texttt{game\_result = LOSE} \\
            }
        }
    }
    & & \\ \\

\end{supertabular}
}

\end{document}
